\section{HexCore Phases 45--47: Grounded Evidence, Continual Learning, and a Replaceable Language Substrate}
\label{sec:hexcore-phases-45-47-grounding-continual-substrate}

\subsection{Programme objective}

Phases 45--47 connect three previously separate requirements of the AION
cognitive architecture:
\begin{enumerate}
    \item grounding decisions in provenance-bearing evidence from several
    modalities;
    \item accumulating verified procedural experience without forgetting
    earlier capabilities; and
    \item using a substantially stronger pretrained language model for rapid
    semantic proposals without transferring truth or promotion authority to
    that model.
\end{enumerate}

The combined architecture follows the invariant
\[
\begin{aligned}
\text{neural or perceptual proposal}
&\rightarrow \text{HexCore compatibility check}\\
&\rightarrow \text{executable verification}\\
&\rightarrow \text{accept, abstain, or fall back}\\
&\rightarrow \text{outcome memory}
\rightarrow \text{CAU}.
\end{aligned}
\]
Here, Cognitive Authority Unification (CAU) remains the final promotion
authority.  Language models, modality parsers, learned routers, and retrieved
procedures remain proposal mechanisms.

\subsection{Phase 45: multimodal world grounding with GlyphChain proofs}
\label{subsec:phase45-multimodal-glyphchain}

\subsubsection{Research question}

Phase 45 tested whether AION could ground a decision in several distinct
evidence modalities, preserve exact provenance, detect contradictions and
bind the resulting evidence set into a tamper-detecting commitment.
GlyphChain was used as a \emph{proof rail}, not as a payment rail.

\subsubsection{Evidence architecture}

Each evaluation world contained four actual artifacts:
\begin{enumerate}
    \item a natural-language policy document containing a limit, revision and
    reported observation;
    \item a CSV table containing the primary measured value;
    \item a PGM pixel image containing a visually encoded gauge reading; and
    \item a JSON software-state record containing the active state and current
    policy revision.
\end{enumerate}

Every modality-specific parser produced an
\texttt{aion.multimodal.evidence\_capsule.v1} capsule containing:
\begin{itemize}
    \item a stable evidence identifier and modality;
    \item source URI and exact evidence location;
    \item a content hash and observation time;
    \item the extracted value and extraction confidence; and
    \item whether the value was directly parsed or merely reported.
\end{itemize}

The governed decision policy accepted a case only when the image and table
readings agreed, the reported reading agreed with the table, the policy
revision matched the live software revision, the software control state was
active, and the verified reading did not exceed the policy limit.  A verified
reading above the limit produced a rejection.  Any unresolved cross-modal
conflict produced abstention rather than a guessed decision.

\subsubsection{GlyphChain evidence commitment}

The final decision, contradiction set and canonical evidence-capsule hashes
were inserted into an \texttt{AION\_EVIDENCE\_PROOF\_V1} envelope.  The
GlyphChain canonical codec generated the commitment hash.  Verification
recomputed both the payload hash and the commitment hash:
\[
C_{\mathrm{evidence}}
=
H\!\left(
  \operatorname{canon}
  \left[
    D,\,
    \mathcal{C},\,
    H(E_1),\ldots,H(E_n)
  \right]
\right),
\]
where \(D\) is the governed decision, \(\mathcal{C}\) is the contradiction
set, and \(E_i\) is a provenance-bearing evidence capsule.  Each evaluation
also mutated the committed decision after hashing; every mutated envelope was
required to fail verification.

This adapter generated and verified deterministic internal commitments.  It
did not submit a live GlyphChain transaction, move PHO, access a wallet, or
create a payment or escrow obligation.

\subsubsection{Phase 45 sealed results}

The official evaluation contained 140 worlds: 100 sealed worlds and 40
external-namespace worlds distributed across five contradiction families.

\begin{table}[htbp]
\centering
\begin{tabular}{lr}
\hline
\textbf{Measure} & \textbf{Result}\\
\hline
Multimodal decision accuracy & 100\%\\
Weakest-family accuracy & 100\%\\
External accuracy & 100\%\\
Text-only control accuracy & 40\%\\
Accuracy improvement over text-only control & +60 points\\
Provenance-complete evidence sets & 100\%\\
Verified GlyphChain commitments & 100\%\\
Tampered commitments accepted & 0\\
Unsafe false acceptances & 0\\
Restart retention & 100\%\\
Live chain transactions & 0\\
Payment side effects & 0\\
\hline
\end{tabular}
\caption{Phase 45 multimodal grounding and proof results.}
\label{tab:phase45-multimodal-results}
\end{table}

The promoted procedure was
\texttt{procedure\_multimodal\_grounding\_9bc5141c1e09}.
The dedicated GlyphChain proof suite passed \(9/9\) tests, and the complete
HexCore history passed \(55/55\) tests at the time of promotion.

\subsubsection{Phase 45 boundary}

The image modality was a controlled scalar pixel gauge and each parser was
engineered.  The result therefore establishes bounded cross-modal grounding,
provenance and tamper detection; it does not establish general computer
vision, arbitrary diagram interpretation, audio/video cognition, live sensor
deployment or artificial general intelligence.

\subsection{Phase 46: large continual-learning arena}
\label{subsec:phase46-continual-arena}

\subsubsection{Research question and arena}

Phase 46 tested whether AION could accumulate verified experience across
different reasoning families, use retained experience to select better
procedures, preserve earlier competence while learning new skills, abstain on
unknown operators and retain all improvements after restart.

The official arena contained 2,040 tasks:
\begin{table}[htbp]
\centering
\begin{tabular}{lr}
\hline
\textbf{Cohort} & \textbf{Tasks}\\
\hline
Development & 1,200\\
Sealed & 480\\
External-source phrasing & 240\\
Unknown-operator OOD & 120\\
\hline
Total & 2,040\\
\hline
\end{tabular}
\caption{Phase 46 arena composition.}
\label{tab:phase46-arena-composition}
\end{table}

Twelve executable procedure families covered mathematics, formal logic,
evidence support and contradiction, causal state transitions, sequential and
branching planning, and numeric and Photon tool routing.  A learned router
could propose a procedure, but could not approve it.  Every accepted answer
was checked by an executable authority.  Rejected or low-confidence proposals
fell back to governed exhaustive search.

\subsubsection{Continual challenger protocol}

Training proceeded through three private challenger generations:
\begin{itemize}
    \item Generation 1 learned mathematics and logic.
    \item Generation 2 learned evidence and causal reasoning while replaying
    Generation 1.
    \item Generation 3 learned planning and tool routing while replaying all
    preceding families.
\end{itemize}

Promotion required unchanged mean and weakest-family accuracy, fewer real
procedure attempts, at least 95\% first-choice transfer, at least 98\%
backward retention, zero unsafe accepted proposals, source-disjoint sealed
evaluation and CAU authorization.  Unknown operators had no compatible
executable contract and were therefore required to abstain irrespective of
router confidence.

\subsubsection{Compounding efficiency}

\begin{table}[htbp]
\centering
\begin{tabular}{llrr}
\hline
\textbf{Generation} & \textbf{New capability family}
& \textbf{Cumulative tasks} & \textbf{Attempt reduction}\\
\hline
1 & Mathematics and logic & 400 & 17.72\%\\
2 & Evidence and causal reasoning & 800 & 51.42\%\\
3 & Planning and tool routing & 1,200 & 84.46\%\\
\hline
\end{tabular}
\caption{Phase 46 cumulative continual-learning efficiency.}
\label{tab:phase46-generations}
\end{table}

All generations retained 100\% mean and weakest-family outcome accuracy.  A
control trained only on the latest generation achieved a 17.49\% attempt
reduction, whereas the cumulative learner achieved 84.46\%.  The resulting
continual-memory advantage was therefore
\[
84.46 - 17.49 = 66.97\ \text{percentage points}.
\]

\begin{table}[htbp]
\centering
\begin{tabular}{lr}
\hline
\textbf{Final measure} & \textbf{Result}\\
\hline
Outcome accuracy & 100\%\\
Weakest-family accuracy & 100\%\\
First-choice procedure accuracy & 100\%\\
Reasoning-attempt reduction & 84.46\%\\
Maximum measured forgetting & 0\%\\
OOD safe abstention & 100\%\\
OOD unsafe acceptances & 0\\
Restart retention & 100\%\\
\hline
\end{tabular}
\caption{Phase 46 final governed results.}
\label{tab:phase46-final-results}
\end{table}

The final promoted procedure was
\texttt{procedure\_continual\_arena\_g3\_ed42d407dc17}.  At completion,
the complete HexCore history passed \(56/56\) tests.

\subsubsection{Phase 46 interpretation and boundary}

The experiment establishes bounded continual procedural learning: earlier
verified outcomes materially changed which reasoning method AION tried first,
later learning did not erase earlier tested competence, unfamiliar operations
triggered safe fallback, and the learned state survived restart.  The arena
was procedurally generated rather than a collection of 2,040 independently
human-authored problems.  It therefore does not establish unrestricted
domain-general reasoning, broad social or creative intelligence, or AGI.

\subsection{Phase 47: stronger replaceable cognitive substrate}
\label{subsec:phase47-gemma-substrate}

\subsubsection{Objective and substrate contract}

Phase 47 tested whether a substantially stronger pretrained language
substrate could improve natural-language transfer without acquiring authority
over truth, memory, execution or promotion.

The local frozen substrate had the following recorded configuration:
\begin{table}[htbp]
\centering
\begin{tabular}{ll}
\hline
\textbf{Property} & \textbf{Value}\\
\hline
Runtime model & \texttt{gemma4:e2b}\\
Parameters & 5.1 billion\\
Quantization & Q4\_K\_M\\
Context window & 131,072 tokens\\
Representation width & 1,536\\
Local artifact size & approximately 7.2 GB\\
Advertised capabilities & completion, vision, audio, tools, thinking\\
\hline
\end{tabular}
\caption{Phase 47 local substrate configuration.}
\label{tab:phase47-substrate}
\end{table}

The recorded model digest was:
\begin{center}
\small\ttfamily
7fbdbf8f5e45a75bb122155ed546e765\\
b4d9c53a1285f62fd9f506baa1c5a47e
\end{center}
The installed Ollama package did not expose this model through its embedding
endpoint.  The implementation therefore used deterministic constrained
generation with a strict procedure-label schema and a small output budget.
Responses were cached by model and instruction hash, allowing interrupted
evaluation to resume safely.

\subsubsection{Authority-preserving architecture}

\[
\begin{aligned}
\text{natural instruction}
&\rightarrow \text{frozen Gemma proposal}\\
&\rightarrow \text{strict label parser}\\
&\rightarrow \text{HexCore capability-contract check}\\
&\rightarrow \text{executable verification}\\
&\rightarrow \text{accept or governed fallback}\\
&\rightarrow \text{outcome memory and CAU}.
\end{aligned}
\]

The substrate was explicitly prohibited from accepting its own proposal,
creating an executable compatibility contract, writing HexCore knowledge,
bypassing fallback, approving promotion, modifying the protected 150M
foundation, or becoming necessary for correct governed operation.

\subsubsection{Evaluation contract}

The evaluation contained 96 sealed natural-paraphrase tasks, 48
external-source paraphrase tasks and 48 subjective, ambiguous or unknown OOD
requests.  Twelve procedure classes covered mathematics, logic, evidence,
causal reasoning, planning and tool routing.  A lexical control learned only
from the procedural language of Phase 46.  Gemma received the same allowed
procedure schema but source-disjoint natural instructions.

\subsubsection{Phase 47 results}

\begin{table}[htbp]
\centering
\begin{tabular}{lr}
\hline
\textbf{Measure} & \textbf{Result}\\
\hline
Lexical-control proposal accuracy & 8.33\%\\
Gemma proposal accuracy & 90.97\%\\
Natural-language transfer gain & +82.64 points\\
Weakest-family proposal accuracy & 66.67\%\\
Governed outcome accuracy & 100\%\\
Reasoning-attempt reduction & 73.66\%\\
Mean recorded local proposal latency & 321 ms\\
Raw Gemma OOD abstention & 50\%\\
Governed OOD abstention & 100\%\\
Provider-disabled governed accuracy & 100\%\\
Unsafe accepted proposals & 0\\
Restart retention & 100\%\\
\hline
\end{tabular}
\caption{Phase 47 replaceable-substrate results.}
\label{tab:phase47-results}
\end{table}

The promoted procedure was
\texttt{procedure\_stronger\_substrate\_3ddbe4956a12}.  The complete HexCore
history passed \(57/57\) tests after integration.

\subsubsection{Decisive negative result}

Raw Gemma abstained correctly on only 50\% of unsupported OOD requests.  This
was retained as an explicit substrate limitation rather than hidden by the
governed score.  The initial evaluation correctly rejected any design that
treated raw model abstention as the safety authority.

The final architecture instead enforced the pre-declared compatibility rule:
a proposal cannot be accepted without a valid executable capability contract.
Every OOD request lacked such a contract and was routed to abstention or
fallback.  Consequently, governed OOD abstention reached 100\%, provider-off
accuracy remained 100\%, and no unsafe proposal was accepted.

This result demonstrates the intended dual-process separation:
\[
\underbrace{\text{pretrained neural substrate}}_{\text{semantic intuition}}
\quad+\quad
\underbrace{\text{HexCore and executable authorities}}_{\text{verification,
memory, and control}}.
\]
The neural model improves semantic proposal quality; it does not determine
what AION is permitted to accept as knowledge or action.

\subsection{Combined architectural achievement}

Together, Phases 45--47 establish the following bounded cognitive chain:
\[
\begin{aligned}
\text{multimodal artifacts}
&\rightarrow \text{provenance-bearing evidence capsules}\\
&\rightarrow \text{contradiction-aware grounded decision}\\
&\rightarrow \text{GlyphChain tamper-evident commitment}\\
&\rightarrow \text{verified outcome memory}\\
&\rightarrow \text{continually improved procedure selection}\\
&\rightarrow \text{stronger natural-language proposals}\\
&\rightarrow \text{HexCore verification and CAU}.
\end{aligned}
\]

The central advance is architectural rather than a claim of frontier
intelligence.  AION can combine grounded evidence, persistent experience and
a strong local semantic model while remaining operational when that model is
removed.  Its accepted decisions retain evidence, execution and governance
links instead of depending solely on an opaque language-model output.

\subsection{Recorded limitations and next evaluation}

The experiments remain bounded:
\begin{itemize}
    \item multimodal images were controlled gauges rather than unrestricted
    natural scenes;
    \item the continual arena used engineered procedural tasks and authorities;
    \item the language benchmark used constrained internal procedure
    classification rather than broad discourse or generation;
    \item the local Gemma substrate was not fine-tuned and did not become an
    AION-native language generator;
    \item no result establishes frontier competitiveness, unrestricted
    autonomous learning or AGI.
\end{itemize}

The next numbered stage is Phase 48: frozen, contamination-controlled external
evaluation.  It must use independently sourced human-authored tasks, fixed
protocols, source and training isolation, public baselines, appropriate
executable or human authorities, and third-party reproducibility.  Phase 47
V2 should separately test the frozen substrate on natural-document reading,
open-schema extraction, multimodal proposal generation and grounded language
generation while preserving the same authority boundary.
